
%%=============================================================%%
%% DISCUSSION                                                  %%
%%=============================================================%%
\section{Discussion}
\label{sec:discussion}

Deep Conformal Alignment shows that morphological variability in histopathology,
a principal driver of poor data efficiency and limited cross-institution
transfer, can be addressed by learning to normalise tissue geometry as part of
the classification pipeline rather than as a fixed preprocessing step. Three
findings are of practical significance. First, the accuracy DCA gains for a given
tissue class scales with that class's morphological heterogeneity
($r=0.94$), so the method concentrates its benefit exactly where irregularity is
greatest---tumour, stroma, and debris---and reduces confusion between
morphologically overlapping classes. Second, DCA is markedly label-efficient,
matching supervised baselines that use roughly three-to-four times more
annotations and outperforming self-supervised pretraining in the low-data regime;
this matters most where expert labels are scarce, such as rare subtypes and
resource-limited settings. Third, geometric normalisation is complementary to
foundation-model pretraining: adding DCA's normalisation head to a pretrained
encoder improves both accuracy and cross-institution robustness beyond either
component alone, so the approach remains useful in the foundation-model era rather
than being superseded by it.

The learned transformations are also interpretable. Their magnitude adapts to
tissue type, they remain approximately angle-preserving (92\% of pixels below a
low conformal-energy threshold), and the residual conformal energy is elevated
precisely on the patches the model misclassifies, pointing to conformal energy as
a candidate intrinsic uncertainty measure for clinical decision support.

Several limitations define the scope of these conclusions and motivate future
work. Both datasets are colorectal cancer with H\&E staining; generalisation to
other organs (breast, lung, brain) and to other modalities (immunohistochemistry,
fluorescence) remains to be established, although the underlying geometric
principles are not tissue-specific. Extension to gigapixel whole-slide analysis
would require hierarchical architectures that maintain geometric coherence across
scales. Establishing a calibrated relationship between conformal energy and
predictive confidence could turn the interpretability signal observed here into a
usable uncertainty estimate. Finally, integrating geometric normalisation with
multimodal foundation models that combine image, text, and molecular data is a
natural next step toward more comprehensive diagnostic systems. Overall, DCA
indicates that building domain-appropriate geometric structure into deep learning
is a productive route to more label-efficient, robust, and interpretable clinical
image analysis.
